%! Carrying Out a Test for the Difference of Two Population Means — Unit 7, Lesson 7.9 — Lesson Plan
\documentclass[10pt]{article}
\usepackage{apstats-article}
\usepackage{apstats-boxes}
\usepackage{graphicx}
\graphicspath{{images/}}
\newcommand{\TallMath}[1]{$\displaystyle #1\rule[-1.4em]{0pt}{3.2em}$}
\newcommand{\UnitNumberName}{Unit 7: Inference for Quantitative Data: Means \quad}
\newcommand{\LessonNumberName}{Lesson 7.9: Carrying Out a Test for the Difference of Two Population Means}

\begin{document}
\begin{center}
    {\Large\bfseries \CourseName: \SchoolYear \\
    {\small\bfseries \UnitNumberName \LessonNumberName}}
\end{center}

\begin{tcolorbox}[colback=sky, colframe=navy, boxrule=0.9pt, arc=2mm,
    left=2mm, right=2mm, top=1.5mm, bottom=1.5mm]
    \textbf{Primary Objective:} Students will carry out a two-sample $t$-test for a difference
    of population means by computing the $t$-statistic from summary statistics, finding the
    $p$-value using technology, interpreting the $p$-value assuming $H_0$ is true, and making
    a formal, context-grounded conclusion by comparing $p$ to $\alpha$.
    (Big Ideas: VAR, DAT)
\end{tcolorbox}

\renewcommand{\arraystretch}{1.6}
\begin{skillbox}[Priority Ideas \& Skills]{goldbox}
\begin{minipage}[t]{0.48\linewidth}
    \textbf{AP Skill 3 --- Using Probability \& Simulation}
    \begin{itemize}
        \item Construct the two-sample $t$-statistic from the general formula and the
              SE formula:
              $t = \dfrac{(\bar x_1 - \bar x_2) - 0}{\sqrt{s_1^2/n_1 + s_2^2/n_2}}$
              (Skill 3.E; VAR-7.I)
        \item Use technology to obtain the $p$-value and degrees of freedom for a
              two-sample $t$-test.
    \end{itemize}
    \textbf{AP Skill 4 --- Statistical Argumentation}
    \begin{itemize}
        \item Interpret the $p$-value of a two-sample $t$-test in context, recognizing
              it is computed assuming $H_0: \mu_1 - \mu_2 = 0$ (Skill 4.B; DAT-3.G).
        \item Justify a claim about two populations by comparing $p$ to $\alpha$ and
              writing a conclusion in context (Skill 4.E; DAT-3.H).
    \end{itemize}
\end{minipage}
\hfill
\begin{minipage}[t]{0.48\linewidth}
    \textbf{Key Understandings}
    \begin{itemize}
        \item The two-sample $t$-statistic follows the same general structure as all
              test statistics: (estimate $-$ null value) $\div$ SE.
        \item The SE for $\bar x_1 - \bar x_2$ is $\sqrt{s_1^2/n_1 + s_2^2/n_2}$,
              combining the variability from both groups.
        \item Degrees of freedom are found with technology; they fall between
              $\min(n_1 - 1,\, n_2 - 1)$ and $n_1 + n_2 - 2$.
        \item The $p$-value is computed \emph{assuming the null is true} ($\mu_1 = \mu_2$)
              --- it is not the probability that $H_0$ is true.
        \item The conclusion restates the decision in terms of the original research
              question with appropriate uncertainty language.
    \end{itemize}
\end{minipage}
\end{skillbox}

\begin{skillbox}[{Vocabulary, Concepts, \& Theorems}]{greenbox}
\begin{minipage}[t]{0.48\linewidth}
    \begin{itemize}
        \item \textbf{Two-sample $t$-statistic} ---
              \TallMath{t = \dfrac{(\bar x_1 - \bar x_2) - (\mu_1 - \mu_2)_0}
              {\sqrt{s_1^2/n_1 + s_2^2/n_2}}}; under $H_0$, $(\mu_1-\mu_2)_0 = 0$
        \item \textbf{SE of $\bar x_1 - \bar x_2$} ---
              $\sqrt{s_1^2/n_1 + s_2^2/n_2}$; pooled from each group's variance
        \item \textbf{Degrees of freedom} --- calculated by technology; bounded by
              $\min(n_1-1, n_2-1)$ and $n_1+n_2-2$
    \end{itemize}
\end{minipage}
\hfill
\begin{minipage}[t]{0.48\linewidth}
    \begin{itemize}
        \item \textbf{$p$-value (two-sample)} --- probability of observing a
              $t$-statistic as extreme as the computed one, assuming $\mu_1 = \mu_2$
        \item \textbf{Two-sample $t$-test} --- inference procedure for comparing two
              independent group means when $\sigma_1$ and $\sigma_2$ are unknown
        \item \textbf{Formal decision} --- reject $H_0$ if $p \le \alpha$;
              fail to reject $H_0$ if $p > \alpha$; never ``accept'' $H_0$
    \end{itemize}
\end{minipage}
\end{skillbox}

\begin{skillbox}[Activate Prior Knowledge \& Spiral Review]{sky}
\begin{minipage}[t]{0.55\linewidth}\vspace{0pt}
    \textbf{Spiral Review (warm-up)}
    \begin{itemize}
        \item Lesson 7.8: state hypotheses and verify conditions for a two-sample $t$-test
        \item Lesson 7.5: compute a one-sample $t$-statistic and find the $p$-value using
              technology --- the same structure applies here
        \item SE formula from the AP formula sheet:
              $SE_{\bar x_1 - \bar x_2} = \sqrt{s_1^2/n_1 + s_2^2/n_2}$
    \end{itemize}
\end{minipage}
\hfill
\begin{minipage}[t]{0.38\linewidth}\vspace{0pt}\centering
    % Authored warm-up compiles to target/ --- no source PDF to embed here.
\end{minipage}
\end{skillbox}

\begin{skillbox}[Hook]{sky}
\begin{minipage}[t]{0.55\linewidth}
    \textbf{ACL Surgery Recovery (3--4 min)}\\
    Return to the setup from Lesson 7.8: a 2018 AP exam study compared recovery times (days)
    after ACL surgery. Group 1 used a new surgical technique ($n_1 = 110$, $\bar x_1 = 47.2$,
    $s_1 = 8.4$); Group 2 used the standard technique ($n_2 = 100$, $\bar x_2 = 36.7$,
    $s_2 = 9.1$). We already stated $H_0: \mu_1 - \mu_2 = 0$ and verified conditions.\\[4pt]
    Ask: ``We know the two sample means differ by about 10.5 days. Is that gap large enough
    --- relative to the natural variability in the data --- to reject the null?''\\[4pt]
    Reveal: $SE = \sqrt{8.4^2/110 + 9.1^2/100} \approx 1.47$, so
    $t \approx (47.2 - 36.7)/1.47 \approx 7.13$.\\[4pt]
    ``How do we turn $t \approx 7.13$ into a probability we can act on?''
\end{minipage}
\hfill
\begin{minipage}[t]{0.40\linewidth}
    \textbf{What You Already Know}
    \begin{itemize}
        \item Test statistic $=$ (estimate $-$ null) $\div$ SE
              --- same skeleton as all AP tests.
        \item The SE formula for $\bar x_1 - \bar x_2$ is on the AP formula sheet;
              the $t$ formula is not, but can be assembled from these pieces.
        \item The $p$-value always answers: ``If $H_0$ were true, how likely is a
              result this extreme?''
        \item Decision rule ($p \le \alpha \Rightarrow$ reject) is universal.
    \end{itemize}
\end{minipage}
\end{skillbox}

\begin{skillbox}[Lesson: Computing the Two-Sample $t$-Statistic]{sky}
\begin{minipage}[t]{0.48\linewidth}
    \textbf{The General Formula Applied}
    \[
      t = \frac{\text{estimate} - \text{null value}}{SE}
        = \frac{(\bar x_1 - \bar x_2) - 0}{\sqrt{s_1^2/n_1 + s_2^2/n_2}}
    \]
    \begin{itemize}
        \item Numerator: observed difference minus null difference (0)
        \item Denominator: $SE = \sqrt{s_1^2/n_1 + s_2^2/n_2}$ (from formula sheet)
        \item Result follows an approximate $t$-distribution; $df$ from technology
        \item $df$ range: $\min(n_1-1, n_2-1)$ to $n_1+n_2-2$
    \end{itemize}
    \textit{ACL example: $SE \approx 1.47$, $t \approx 7.13$, $df \approx 207$.}
\end{minipage}
\hfill
\begin{minipage}[t]{0.48\linewidth}
    \textbf{Finding the $p$-Value with Technology}
    \begin{itemize}
        \item \textbf{One-sided ($H_a: \mu_1 - \mu_2 > 0$):} area to the right of $t$
        \item \textbf{One-sided ($H_a: \mu_1 - \mu_2 < 0$):} area to the left of $t$
        \item \textbf{Two-sided ($H_a: \mu_1 \ne \mu_2$):} $2 \times$ area beyond $|t|$
    \end{itemize}
    \textbf{TI-84: 2-SampTTest} (STAT $\to$ TESTS $\to$ 4)\\
    Enter $\bar x_1, s_1, n_1, \bar x_2, s_2, n_2$; select $H_a$; \textbf{do not pool}.\\[4pt]
    \textit{ACL ($H_a: \mu_1 > \mu_2$): $p = P(t > 7.13,\ df \approx 207) \approx 0$.}\\[4pt]
    \textbf{Key reminder:} use \texttt{Pooled: No} unless explicitly told to assume equal $\sigma$.
\end{minipage}
\end{skillbox}

\begin{skillbox}[Lesson (cont.): Interpret and Conclude]{sky}
\begin{minipage}[t]{0.48\linewidth}
    \textbf{Interpreting the $p$-Value (DAT-3.G)}\\
    The $p$-value is the probability of observing a difference in sample means as large as
    (or larger than) the one observed, \emph{if the true population means are equal}.\\[6pt]
    \textit{ACL ($p \approx 0$):} ``If the two surgical techniques truly produce the same mean
    recovery time, there is essentially a 0\% chance of seeing sample means this far apart
    by chance alone.''\\[4pt]
    \textbf{Common error:} ``The $p$-value is the chance the groups are equal.''
    --- Wrong. $H_0$ is \emph{assumed} true; $p$ is computed \emph{under} that assumption.
\end{minipage}
\hfill
\begin{minipage}[t]{0.48\linewidth}
    \textbf{Formal Decision and Conclusion (DAT-3.H)}\\
    Compare $p$ to $\alpha = 0.05$:
    \begin{itemize}
        \item $p \approx 0 \le 0.05$ $\Rightarrow$ \textbf{Reject $H_0$}.
    \end{itemize}
    \textbf{Conclusion template:}\\
    ``Because $p \approx 0 < \alpha = 0.05$, we reject $H_0$. We have convincing evidence
    that the true mean recovery time with the new technique is greater than with the standard
    technique.''\\[4pt]
    \textit{Key phrases:} ``convincing evidence that [$H_a$ in context]'' when rejecting;
    ``not sufficient evidence'' when failing to reject.
\end{minipage}
\end{skillbox}

\begin{skillbox}[Active Monitoring]{redbox}
\begin{minipage}[t]{0.48\linewidth}
    \textbf{Circulate and Check}
    \begin{itemize}
        \item Students use $s_1$ and $s_2$ (not a pooled $s$) in the SE formula.
        \item $p$-value direction matches $H_a$ (greater, less, or two-sided).
        \item $p$-value interpretation says ``assuming the null is true'' --- not
              ``the probability that the groups are equal.''
        \item Conclusion is in context, references $\alpha$, and uses ``convincing
              evidence'' or ``not sufficient evidence.''
        \item Students select \textbf{Pooled: No} on the TI-84 2-SampTTest.
    \end{itemize}
\end{minipage}
\hfill
\begin{minipage}[t]{0.48\linewidth}
    \textbf{Cold-Call Prompts}
    \begin{itemize}
        \item ``What does the denominator $\sqrt{s_1^2/n_1 + s_2^2/n_2}$ represent?''
        \item ``The $p$-value is 0.12 and $\alpha = 0.05$. What do you conclude?''
        \item ``How does the two-sample $t$-statistic differ from a one-sample $t$?''
        \item ``If $H_a$ is two-sided, how does the $p$-value calculation change?''
        \item ``Why don't we need to know $\sigma_1$ or $\sigma_2$?''
    \end{itemize}
\end{minipage}
\end{skillbox}

\begin{skillbox}[Group Work \& Differentiation]{redbox}
\begin{multicols}{3}
    \textbf{Tier R --- Remediate}
    \begin{itemize}
        \item Provide a formula strip with $SE = \sqrt{s_1^2/n_1 + s_2^2/n_2}$ and a
              completed TI-84 2-SampTTest example.
        \item Students read a completed $t$-statistic and $p$-value and write the
              interpretation and conclusion sentences.
    \end{itemize}
    \columnbreak
    \textbf{Tier A --- Approaching}
    \begin{itemize}
        \item Given two-group summary statistics, compute $SE$ and $t$ by hand,
              confirm with technology, and write full SPDC conclusion.
        \item Interpret the $p$-value in a complete sentence referencing $H_0$.
    \end{itemize}
    \columnbreak
    \textbf{Tier E --- Extension}
    \begin{itemize}
        \item Compare one-sided vs.\ two-sided $p$-values for the same data; discuss
              which alternative is appropriate and the consequences of choosing wrong.
        \item Explain how increasing sample sizes would affect $SE$, $t$, and the
              decision.
    \end{itemize}
\end{multicols}
\end{skillbox}

\begin{skillbox}[Individual Work \& Assessment]{redbox}
\begin{minipage}[t]{0.48\linewidth}
    \textbf{Exit Ticket (5 min)}\\
    Sleep researchers compared sleep duration for college athletes vs.\ non-athletes.
    Summary: $\bar x_1 = 7.4$ h, $s_1 = 0.9$ h, $n_1 = 30$;
    $\bar x_2 = 6.8$ h, $s_2 = 1.1$ h, $n_2 = 28$. Technology gives $t \approx 2.30$,
    $df \approx 53$, $p \approx 0.025$ (one-sided, right).
    \begin{enumerate}
        \item Interpret the $p$-value in context.
        \item Write a complete conclusion at $\alpha = 0.05$.
    \end{enumerate}
\end{minipage}
\hfill
\begin{minipage}[t]{0.48\linewidth}
    \textbf{AP-Style Check}\\
    \textit{A two-sample $t$-test gives $t = 1.65$, $df = 45$, $p = 0.053$ (one-sided).
    At $\alpha = 0.05$, the correct conclusion is:}
    \begin{enumerate}[label=(\Alph*)]
        \item Reject $H_0$; there is convincing evidence of a difference.
        \item Fail to reject $H_0$; there is not sufficient evidence.
        \item Accept $H_0$; the populations have equal means.
        \item Reject $H_a$; the data support the null.
    \end{enumerate}
    Answer: (B)
\end{minipage}
\end{skillbox}

\begin{skillbox}[Reinforcement \& Extension]{goldbox}
\begin{minipage}[t]{0.48\linewidth}
    \textbf{Homework Overview}
    \begin{itemize}
        \item Four full SPDC problems: two one-sided and two two-sided two-sample $t$-tests
              in varied real-world contexts.
        \item Students compute SE and $t$ by hand, verify with technology, interpret the
              $p$-value, and write formal conclusions.
        \item Reflection: explain what it means to fail to reject $H_0$ when comparing two groups.
    \end{itemize}
\end{minipage}
\hfill
\begin{minipage}[t]{0.48\linewidth}
    \textbf{Extension / Preview}
    \begin{itemize}
        \item \textit{Extension:} Re-run the ACL example with $H_a: \mu_1 \ne \mu_2$ (two-sided).
              Compare the $p$-value and conclusion to the one-sided result.
        \item \textit{Preview:} Lesson 7.10 focuses on selecting the correct inference
              procedure --- when to use one-sample $t$, two-sample $t$, matched pairs, or
              other procedures --- and connecting all of Unit 7 together.
    \end{itemize}
\end{minipage}
\end{skillbox}

\end{document}
